\documentclass[a4paper,10pt]{article}
\usepackage[francais]{babel}
\usepackage{url}

\selectlanguage{francais}
\input{letter-LSIIT-UdS.sty}

\selectlanguage{francais}
\begin{document}


\name{St\'ephane \textsc{Genaud}}
\fonction{Ma\^itre de conf\'erences (HDR)}
\telephoneillkirch{+33 (0)3688 54542}
\email{genaud@unistra.fr}

\begin{letter}[Recommendation of Marc Fr\^incu for position 3121 / EPL 016 / ICTM / 2013]{to}{Universit\'e Catholique de Louvain}

Dear Madam/Sir, 

I first met Marc  Fr\^incu in 2007 when he was a student  on an internship in my
team at the  INRIA research center of Nancy,  France. I was at that  time on the
leave from my university on a researcher position. Although we have not directly
worked  together at  that time,  I  know from  my colleagues  that he  perfectly
achieved  his objectives  during  his  internship, which  led  my colleagues  to
propose him to apply for a \emph{CORDI} PhD funding at INRIA.

Marc preferred to  return to Romania, where he  went on with a PhD  on a subject
related to scheduling applications for Grids. This research theme is close to my
owns, and I  noticed during the last years the  interesting results published by
Marc in conferences in the domain.

During last summer, I learned that Marc was searching for a PostDoc position and
I immediately proposed  him to join us,  to work on scheduling  for clouds. This
theme perfectly matched  Marc's professional objectives. Marc joined  us on last
October,       to      participate       in      the       national      project
SONGS\footnote{http://infra-songs.gforge.inria.fr/}  I  am   involved  in.   The
objective  is to  enrich the  existing SimGrid  simulation toolkit  with an  API
dedicated to the  simulation of cloud-based systems. Simultaneously  to the work
on  simulation, we  carry  on theoretical  studies on  the  impact of  different
scheduling strategies  on different types  of workloads, that may  correspond to
certain  activity  types,  or  certain application  models  (e.g  submission  of
independent  jobs, or  workflows,  etc).  These  scheduling  algorithms will  be
candidate solutions  that could  benefit from  our cloud  simulator development.
Ideally, our  simulation results will  be compared to  those produced in  a real
software framework, such as the one under development in our team.

Since he  arrived, Marc has  actively participated in  all these tasks.   He has
shown an  impressive body of  general knowledge in the  field of grid  and cloud
computing. I  asked him to recap  works from the  past year where he  thought we
could work  on, and he came  up with a very  complete survey after a  very short
time. Among other tasks,  he his currently helping a PhD  student who is writing
his first article  by challenging him on potential weaknesses  of the work. Marc
will  also have  a Master  student  under is  supervision  to help  him work  on
scheduling workflows on IaaS clouds.

As a conclusion, I have been favourably  impressed so far by the capabilities of
Marc to  quickly identify and  summarize relevant  research issues, and  then to
write it in a very fluent style.  He is both a very autonomous, and hard-working
person.  I hope  Marc's stay (1 year)  at or institution will be  long enough to
achieve the numerous  objective I would like him to  undertake.  The position he
is applying  for at  Universit\'e de Louvain,  entitled \textit{Large  Scale and
  Cloud Computing}  seems totally in the  scope of his expertise.   Therefore, I
strongly recommend Marc for this position.


\end{letter}
\begin{flushright}
\includegraphics[width=.30\textwidth]{signgenaud.jpg}
\end{flushright}
%\newpage
%\bibliographystyle{plain}
%\bibliography{julien_blaise_ED}


\end{document}
